\documentclass[11pt]{article}
\usepackage[a4paper,margin=27mm]{geometry}
\usepackage{amsmath,amssymb,mathtools}
\usepackage{fontspec,unicode-math}
\setmainfont{Latin Modern Roman}
\setmathfont{Latin Modern Math}
\usepackage{xurl,hyperref}
\hypersetup{colorlinks=true,urlcolor=blue,linkcolor=blue}
\newcommand{\sref}[2]{\hyperlink{source:#1:#2}{[S#2]}}
\newcommand{\cataloguescope}[1]{\paragraph{Mathematical question.} #1}
\begin{document}
\section{Lengths of Bounded-Rank Linear-Order WQOs}\label{statement}

\subsection{Definitions and mathematical statement}

\paragraph{Shared notation.}$\mathbb N=\{1,2,\ldots\}$, and $\mathbb Z,\mathbb Q,\mathbb R,\mathbb C$ denote the integers, rationals, reals and complex numbers. Any other conventions are specified locally.


A countable linear order is scattered if it contains no copy of $\mathbb Q$. Use the Hausdorff hierarchy $\mathcal H_0$ of finite orders; for $\beta>0$, $\mathcal H_\beta$ consists of finite sums of orders $\sum_{n\in\mathbb Z}L_n$ with each $L_n\in\bigcup_{\gamma<\beta}\mathcal H_\gamma$. The Hausdorff rank is the least $\beta$ with $L\in\mathcal H_\beta$. Identify orders $L,K$ when each order-embeds into the other (equimorphism). For a countable ordinal $\alpha$, let $L_\alpha$ be the resulting partial order of scattered orders of rank less than $\alpha$, ordered by embeddability. A quasi-order is wqo if each infinite sequence has increasing indices $i<j$ with $x_i\le x_j$. For a wqo partial order $W$, define $o(W)$ as the supremum of the ordinal order types of its linear extensions.
\paragraph{Statement recorded in the supplied source.}
Determine the maximal order type $o(L_\alpha)$ for every countable ordinal $\alpha$. \sref{statement}{2}
\subsection{Short English statement}
Compute the length of the embeddability ordering of countable linear orders below each Hausdorff-rank bound, after identifying mutually embeddable orders.
\subsection{Sources}
\begin{description}
\item[\hypertarget{source:statement:1}{S1}] ULAM.AI and the UnsolvedMath Contributors, UnsolvedMath: A Curated Collection of Open Mathematics Problems, catalogue entry. CC BY 4.0. \url{https://huggingface.co/datasets/ulamai/UnsolvedMath}.
\item[\hypertarget{source:statement:2}{S2}] Antonio Montalbán, Open Questions in Reverse Mathematics, author preprint (17 August 2010), Question 22, Section 4.2, p. 9 and Definition 4.2, p. 8. \url{https://math.berkeley.edu/~antonio/papers/questionsRM.pdf}.
\end{description}
\label{end:statement}
\end{document}
